\documentclass[14pt,usepdftitle=false,aspectratio=169]{beamer}
\usepackage{preambule}
\setbeamercolor{structure}{fg=black}
\usepackage{probas}\usepackage[nopar]{bigoperators,analyse}\usepackage{complexes,usuelles}\usepackage{matrices}
\hypersetup{pdftitle={Processus stochastiques -- Chaînes de Markov à temps continu}}
\title{Processus stochastiques\\\emph{Chaînes de Markov à temps continu}}
\author{}
\date{}
\begin{document}
\begin{frame}
    \titlepage
\end{frame}
\slideq{Propriété de Markov faible}{1/23}
\slider{Pour tout $t\geqslant0$, conditionnellement à $\calF_t$, $\l X_{t+s}\r_{s\geqslant0}$ est une chaîne de Markov de matrice $Q$ et commençant en $X_t$}{1/23}
\slideq{$x\in E$ est récurrent\linebreak $x\in E$ est transient}{2/23}
\slider{$x$ est récurrent si $\set{t\geqslant0,X_t=x}$ est non borné\linebreak$x$ est transient sinon}{2/23}
\slideq{Théorème de convergence de $P_t$ pour une chaîne de Markov}{3/23}
\slider{Si $X$ est une chaîne de Markov irréductible récurrente positive, et $\pi$ est la mesure de probabilités invariante alors pour tout $x,y\in E$, $P_t\l x,y\r\xrightarrow[t\to+\infty]{}\pi\l y\r$\linebreak En particulier, $\p[\nu]{X_t=y}\xrightarrow[t\to+\infty]{}\pi\l y\r$}{3/23}
\slideq{CNS pour avoir une mesure invariante pour une chaîne irréductible récurrente}{4/23}
\slider{$\mu Q=0$ ssi $\mu P_s=\mu$}{4/23}
\slideq{CNS de non explosion sur les vecteurs propres de $Q$}{5/23}
\slider{$T_\infty=+\infty$ ss'il existe $\theta>0$ tel que, pour tout $v\in\bbR^E$ avec $\left|v_x\right|\leqslant1$ pour tout $x\in E$ et $Qv=\theta v$, alors $v=0$}{5/23}
\slideq{Théorème ergodique}{6/23}
\slider{Si $X$ est une chaîne de Markov irréductible récurrente positive de mesure de probabilités $\pi$, alors $\frac1t\int[s][0][t]{\mathbb1_{X_s=x}}\xrightarrow[t\to+\infty]{}\mu\l x\r$ p.s. pour tout $x$\linebreak Si $f\!:\!E\to\bbR$ est positive bornée $\frac1t\int[s][0][t]{f\l X_s\r}\xrightarrow[t\to+\infty]{}\sum{x\in E}{}{\mu\l x\r f\l x\r}$ p.s.}{6/23}
\slideq{Générateur d'une chaîne de Markov}{7/23}
\slider{$\appl{Q}{E\times E}{\bbR}{\l x,y\r}{\tcase{\lambda\l x\r P\l x,y\r\&\text{si }x\neq y\\\lambda\l x\r\l1-P\l x,x\r\r\&\text{si }x=y\\}}$}{7/23}
\slideq{Propriétés de $Z_x=\esp[x]{\e^{-\theta T_\infty}}$}{8/23}
\slider{Pour tout $x\in E$ et tout $\theta>0$, $\left|Z_x\right|\leqslant1$ et $\theta Z=QZ$\linebreak Si $Z'$ vérifie ces deux conditions alors $Z'\leqslant Z$ terme à terme}{8/23}
\slideq{Critère de non explosion}{9/23}
\slider{{\togglebigopdisplay Pour tout $\theta\ge0$, $\esp{\e^{-\theta T_\infty}}=\esp{\prod{k=0}{+\infty}{\l1+\frac\theta{\lambda\l Z_k\r}\r^{-1}}}$\linebreak$\p{T_\infty=+\infty}=\p{\sum{k=0}{+\infty}{\frac1{\lambda\l Z_k\r}}=+\infty}$\linebreak Dans le cas particulier où $\lambda$ est borné, ou si la loi de $Z_0$ ne charge pas les états transients, alors $T_\infty=+\infty$ p.s.}}{9/23}
\slideq{Équation de Kolmogorov backward}{10/23}
\slider{Si $X$ est un processus càdlàg à valeurs dans $E$, alors $X$ est une chaîne de Markov à temps continu associée à $Q$ si et seulement si pour tout $n\in\bbN$, tous $0\leqslant t_0\leqslant\cdots\leqslant t_{n+1}$ et tous $x_0,\cdots,x_n\in E$, $\p{X_{t_{n+1}}=y\sq X_n=x_n,\cdots,X_0=x_0}=p_{x_n,x_{n+1}}\l t_{n+1}-t_n\r$ où $\l p_{x,y}\r_{x,y}=P$ est la solution minimale de $P\l0\r=I_n$ et $P'\l t\r=QP\l t\r$}{10/23}
\slideq{Chaîne de Markov à temps continu associée à un noyau $P$ et $\lambda\!:\!E\to\bbR_+$}{11/23}
\slider{$\l X_t\r_{t\in\lc0,T_\infty\right[}=\l Z_{N_t}\r$ avec $\l Z_n\r$ une chaîne de Markov sur un ensemble au plus dénombrable $E$ de noyau $P$, $\l E_n^z\r_{n\in\bbN,z\in E}$ des variables exponentielles indépendantes d'intensité $\lambda\l z\r$, $T_0=0$ et $T_n=E_1^{Z_0}+\cdots+E_n^{z_{n+1}}$, $N_t=\sup{u\in\bbN,T_u\leqslant t}\in\bar\bbN$ et $T_\infty=\sum{n=0}{+\infty}{E_{n+1}^{z_n}}$}{11/23}
\slideq{Lien entre $Q$, $T_\cdot$ et la récurrence}{12/23}
\slider{Si $Q\l x,x\r=0$ ou $\p[\!\!x]{T_x<+\infty}=1$ alors $x$ est récurrent et $\int[t][0][+\infty]{P_t\l x,x\r}=+\infty$}{12/23}
\slideq{Équation de Kolmogorov backward}{13/23}
\slider{Si $X$ est un processus càdlàg à valeurs dans $E$, alors $X$ est une chaîne de Markov à temps continu associée à $Q$ si et seulement si pour tout $n\in\bbN$, tous $0\leqslant t_0\leqslant\cdots\leqslant t_{n+1}$ et tous $x_0,\cdots,x_n\in E$, $\p{X_{t_{n+1}}=y\sq X_n=x_n,\cdots,X_0=x_0}=p_{x_n,x_{n+1}}\l t_{n+1}-t_n\r$ où $\l p_{x,y}\r_{x,y}=P$ est la solution minimale de $P\l0\r=I_n$ et $P'\l t\r=P\l t\r Q$}{13/23}
\slideq{Lien entre la récurrence de $x$ dans $X$ est une chaîne de Markov discrète extraite}{14/23}
\slider{Pour $h>0$, $x$ est récurrent pour $\l X_t\r_t$ si et seulement si $x$ est récurrent pour $\l X_{hn}\r_n$}{14/23}
\slideq{Convergence d'une somme de variables aléatoires exponentielles indépendantes}{15/23}
\slider{Si $E_n\sim\expon{\lambda_n}$ alors $\p{\sum{i=0}{+\infty}{E_i}=+\infty}=\mathbb1_{\togglebigopdisplay\sum{k=0}{+\infty}{\frac1{\lambda_k}}}$}{15/23}
\slideq{Propriétés de $\l S_n\r$\linebreak$S_0=0$, $S_{n+1}=\inf{\set{T_\infty>t>S_n,X_t\neq X_{S_n}}}$}{16/23}
\slider{$\l S_n\r$ est strictement croissante et $S_n\to T_\infty$\linebreak Conditionnellement à $S_n<+\infty$ et $X_{S_n}=x$ (ou $\calF_{S_n}$), si $Q\l x,x\r=0$ alors $S_{n+1}=+\infty$ et $x$ est absorbant, et si $Q\l x,x\r\neq0$, alors $S_{n+1}<+\infty$ p.s., $S_{n+1}-S_n\indep X_{S_{n+1}}$, et conditionnellemet à $X_{S_n}=x$, $S_{n+1}-S_n\sim\expon{-Q\l x,x\r}$ et $\p{X_{S_{n+1}}=y\sq X_{S_n}=x}=\oldfrac{-Q\l x,y\r}{Q\l x,x\r}$}{16/23}
\slideq{$\mu$ est une mesure invariante pour $X$}{17/23}
\slider{$\mu$ est non nulle et $\mu Q=0$}{17/23}
\slideq{Fonction de Green de $X$, et lien avec celle sur $Z$}{18/23}
\slider{$G\l x,y\r=\esp[x]{\int[t][0][+\infty]{\mathbb1_{X_t=y}}}$\linebreak$G\l x,y\r=\frac1{\lambda\l y\r}G_Z\l x,y\r$}{18/23}
\slideq{Poissonisation d'une chaîne de Markov}{19/23}
\slider{$\l Z_{N_t}\r_{t\geqslant0}$ pour $\l Z_n\r_{n\in\bbN}$ une chaîne de Markov sur un ensemble au plus dénombrable $E$ et $\l N_t\r_{t\geqslant0}$ un processus de Poisson d'intensité $\lambda$ indépendant de la chaîne de Markov\linebreak$\p{Z_{N_t}=y\sq Z_0=x}=\e^{-\lambda t}\sum{n=0}{+\infty}{\frac{\l\lambda t\r^n}{n!}P^n\l x,y\r}$ pour $P$ le noyau de la chaîne de Markov}{19/23}
\slideq{Processus de Poisson composé}{20/23}
\slider{$\l Z_{N_t}\r_{t\geqslant0}$ pour $\l Z_n\r_{n\in\bbN}$ une marche aléatoire sur un groupe $G$ et $\l N_t\r_{t\geqslant0}$ un processus de Poisson d'intensité $\lambda$ indépendant de la marche aléatoire\linebreak$\p{Z_{N_t}=y\sq Z_0=x}=\e^{-\lambda t}\sum{n=0}{+\infty}{\frac{\l\lambda t\r^n}{n!}P^n\l0,yx^{-1}\r}$ pour $P$ le noyau de la marche aléatoire}{20/23}
\slideq{Mesure invariante en $x$ récurrent}{21/23}
\slider{$\mu^x\l y\r=\esp[x]{\int[s][0][{{T_x^+}}]{\mathbb1_{X_s=y}}}$\linebreak$\mu^x$ est une mesure invariante de masse totale $\esp[x]{T_x^+}$}{21/23}
\slideq{Propriété de Markov forte}{22/23}
\slider{Pour tout temps d'arrêt $T$ pour la filtration canonique, conditionnellement à $\calF_T$, $\l X_{T+s}\r_{s\geqslant0}$ est une chaîne de Markov de matrice $Q$ et commençant en $X_T$}{22/23}
\slideq{Théorème d'existence de mesures invariantes}{23/23}
\slider{Si $X$ admet un état récurrent, alors il existe une mesure invariante\linebreak Si $X$ est irréductible récurrente, toutes les mesures sont proportionnelles}{23/23}
\end{document}